% Part: incompleteness
% Chapter: arithmetization-syntax
% Section: coding-formulas

\documentclass[../../../include/open-logic-section]{subfiles}

\begin{document}

\olfileid{inc}{art}{frm}
\olsection{编码公式}

一旦原始递归地定义了关系 $\fn{Term}(x)$，我们便可用它来原始递归地定义公式对应的关系 $\fn{Frm}(x)$。

\begin{prop}
关系 $\fn{Atom}(x)$（当且仅当 $x$ 是原子公式的Gödel数时成立）是原始递归的。
\end{prop}

\begin{proof}
$x$ 是原子公式的Gödel数，当且仅当以下条件之一成立：
\begin{enumerate}
\item 存在 $n$、$j < x$ 和 $z < x$，使得对每个 $i < n$，
  $\fn{Term}((z)_i)$ 且 $x = $
\[
\Gn{\Obj P^n_j(} \concat \fn{flatten}(z) \concat \Gn{)}.
\]
\item 存在 $z_1, z_2 < x$，使得 $\fn{Term}(z_1)$ 且
  $\fn{Term}(z_2)$，且 $x = $
\[
\Gn{{\eq}(} \concat z_1 \concat \Gn{,} \concat z_2 \concat{\Gn{)}}.
\]
  \tagitem{prvFalse}{$x = \Gn{\lfalse}$。}{}
  \tagitem{prvTrue}{$x = \Gn{\ltrue}$。}{}
\end{enumerate}
\end{proof}

\begin{prop}
\ollabel{prop:frm-primrec}
关系 $\fn{Frm}(x)$（当且仅当 $x$ 是公式的Gödel数时成立）是原始递归的。
\end{prop}

\begin{proof}
符号序列~$s$ 是公式，当且仅当存在一个公式的形成序列~$s_0$, \dots, $s_{k-1} = s$，该序列记录了~$s$ 是如何根据形成规则由原子公式形成的。每个~$s_i$ 的编码（事实上整个序列 $\tuple{s_0, \dots, s_{k-1}}$ 的编码亦然）都小于~$s$ 的编码~$x$。
\end{proof}

\begin{prob}
请仿照 \olref[inc][art][trm]{prop:term-primrec} 的第一个证明的思路，详细给出 \olref[inc][art][frm]{prop:frm-primrec} 的证明。
\end{prob}

\begin{prop}
\ollabel{prop:freeocc-primrec}
关系 $\fn{FreeOcc}(x, z, i)$（当且仅当Gödel数为~$x$ 的公式的第~$i$ 个符号是Gödel数为~$z$ 的变元的自由出现时成立）是原始递归的。
\end{prop}

\begin{proof}
留作练习。
\end{proof}

\begin{prob}
证明 \olref[inc][art][frm]{prop:freeocc-primrec}。你可以利用以下事实：公式中任何本身也是公式的子串，都是该公式的子公式。
\end{prob}

\begin{prop}
性质 $\fn{Sent}(x)$（当且仅当~$x$ 是语句的Gödel数时成立）是原始递归的。
\end{prop}

\begin{proof}
语句是不含变元自由出现的公式。因此 $\fn{Sent}(x)$ 成立当且仅当
\begin{multline*}
\bforall{i<\len{x}}{\bforall{z<x}{}}\\
(\bexists{j<z}{z=\Gn{\Obj v_j}} \lif \lnot\fn{FreeOcc}(x,z,i)).
\end{multline*}
\end{proof}


\end{document}
